\documentclass[10pt,a4paper]{amsart}
\usepackage[margin=19mm]{geometry}
\usepackage{amsmath,amssymb,mathtools}
\usepackage[scr=rsfs,cal=euler]{mathalfa}
\usepackage{xfrac}
\usepackage[unicode,hidelinks,bookmarks=false]{hyperref}
\newcommand{\ie}{i.e.\ }
\newcommand{\cf}{cf.\ }
\newcommand{\resp}{respectively\ }
\newcommand{\Subsection}{\subsection}
\providecommand{\nobreakdash}{\nobreak\hbox{-}}
\setlength{\emergencystretch}{2em}
\multlinegap0pt
\usepackage{enumerate}
\usepackage{amssymb}
\newtheorem{lemma}{Lemma}
\renewcommand{\a}{\alpha}
\renewcommand{\b}{\beta}
\newcommand{\g}{\gamma}
\newcommand{\w}{\omega}
\title{{\bf Relative uniform convergence and Archimedean property in pre-ordered vector spaces and Abelian groups}}
\author{Ivan Emelyanenkov, Eduard Emelyanov}
\date{Source: arXiv:2602.16419v9 (CC BY 4.0)}
\begin{document}
\maketitle

\noindent\emph{Source and licence.} {\bf Relative uniform convergence and Archimedean property in pre-ordered vector spaces and Abelian groups}. Version arXiv:2602.16419v9 (CC BY 4.0), distributed by arXiv under the Creative Commons Attribution 4.0 International licence, http://creativecommons.org/licenses/by/4.0/, as recorded in the arXiv licence metadata for this exact version and shown on the abstract page https://arxiv.org/abs/2602.16419v9. Full licence notice: \texttt{licenses/arxiv-2602.16419-CC-BY-4.0.txt}.\par\smallskip

\noindent\textit{Editorial scope.} Counted unit: the article's lemma H\_is\_closed (source0.tex lines 932--934). The article's complete original proof of it is delivered with it (source0.tex lines 936--986, one \texttt{proof} environment of 51 lines). No mathematics is written, repaired or re-derived: the statement and the proof are the article's own lines, in the article's own notation, and no retained line is substituted or re-flowed.  Retained uncounted as the source-local context the proof needs: lines 915--917.  Nothing was omitted from any retained span, so the package carries no omission record.  No retained line contains a citation, so the wrapper carries no bibliography and no bibliography entry.  No retained line contains a figure environment, an image-inclusion command or drawing code, so no image is delivered and the package declares no asset dependency.  Editorial formatting only, in addition: an AMS \texttt{amsart} preamble that restates the article's own declarations and macros used by the retained lines, namely the environments \texttt{lemma} on separate counters, together with the standard packages those lines need.  The retained environments print in this document as Lemma 1 in order of appearance, whereas the article prints the same items with the numbering of its own full document.  The complete native source of the article as distributed in the arXiv e-print for this exact version is bundled unchanged as \texttt{sources/194880/source0.tex}.
\begin{lemma}\label{there_is_gap}
If $ \b < \a $, then $| 1_\a - x | \not\leqslant t 1_\a$ for any $ x \in H_\a^\b $ and any real number $ 0 \leqslant t < 1 $.
\end{lemma}

\begin{lemma}\label{H_is_closed}
$ (H_\a^{<\b})_{ru}^{(1)} = H_\a^\b $ for every $ 1 < \b < \w_1 $ and $ 1 \leqslant \a < \w_1 $.
\end{lemma}

\begin{proof}
We prove this by induction on $ \alpha $.

\medskip\noindent
\underline{Case $ \a < \b $:}\\
It is obviously true since there is $ \a \leqslant \g < \b $, and thus $ H_\a^\g = H_\a^\b = S_\a $.

\medskip\noindent
\underline{Case $ \a = \b $:}\\
Clearly, $ H_\a^\b = S_\a $, so $ \subseteq $ is trivial.
On the other hand, $ \bigoplus S_{\a(n)} \subseteq H_\a^{<\a} $. Indeed, if we take $ x \in S_\a $ such that $ x_n = 0 $ for all $ n > m $, then, after noticing that
$$
	H_{\a(n)}^{\a(m)} = S_{\a(n)} \quad\text{ for every } n \leqslant m,
$$
we may conclude that $ x \in \bigoplus H_{\a(n)}^{\a(m)} = H_\a^{\a(m)} $.
Now, if $ x \in S_\a $ and $ x_{(m)} $ denotes its initial segment of length $ m $, then $ x_{(m)}\overset{ru}{\rightarrow} x (w) $ 
where the witness $ w $ is given by any positive polynomial growing faster than $ p(x) $, for example, as $ w = (1+t^{\deg p(x) + 1})1_\a $.

\medskip\noindent
\underline{Case $ \a > \b $:}\\
Let's prove $ \supseteq $ first.
For any element $ x \in H_\a^\b $, we know that $ p(x) = 0 $ and $ x_n \in H_{\a(n)}^\b $ for all $ n $, so let $ m $ be any number such that $ x_n = 0 $ for all $ n > m $.
For every $ x_n $ we have sequences of elements $ y_{n, k} \in H_{\a(n)}^{<\b} $ and witnesses 
$ w_n \in S_{\a(n)} $ such that $ y_{n, k}  \overset{ru}{\rightarrow} x_n (w_n) $.
Choose $ \g_k < \b $ such that $ y_{n, k} \in H_{\a(n)}^{\g_k} $ for all $ n < m $, and construct $ y_k \in H_\a^{\g_k} $
$$
	(y_k)_n = 
	\begin{cases}
		y_{n, k}, & \text{ if } n < m, \\
		0, & \text{ if } n \geqslant m,
	\end{cases}
$$
as well as $ w \in S_\a $ from $ w_n $, similarly.
It is easy to see that $ y_k  \overset{ru}{\rightarrow} x (w) $.
\\
Now in the $ \subseteq $ direction.
Let's take an arbitrary sequence $ y_k \in H_\a^{<\b} $ and $ x, w \in S_\a $ such that $ y_k  \overset{ru}{\rightarrow} x (w) $.
Suppose that $ p(x) = p \neq 0 $. Then, for any sufficiently large $ n $, we have $ p(n) \neq 0 $, $ \a(n) \geqslant \b $, $ x_n = p(n) 1_{\a(n)} $, 
and $ w_n = q(n) 1_{\a(n)} $. The ru-convergence gives
$$
	k \ |p(n) 1_{\a(n)} - (y_k)_n | \leqslant q(n) 1_{\a(n)},
$$
and thus
$$
	| 1_{\a(n)} - \frac{1}{p(n)}(y_k)_n | \leqslant \frac{q(n)}{k|p(n)|} 1_{\a(n)}.
$$
But $ y_k \in H_{\a}^\g $ for some $ \g < \b $, which means $ (y_k)_n \in H_{\a(n)}^\g $, which, for large enough $ k $, such that $  \frac{q(n)}{k|p(n)|} < 1 $ 
contradicts Lemma \ref{there_is_gap}.
Hence, $ p(x) = 0 $. On the other hand, $ y_{n, k}  \overset{ru}{\rightarrow} x_n (w_n) $, so, by induction, $ x_n \in H_{\a(n)}^\b $, 
which means that $ x \in \bigoplus H_{\a(n)}^\b = H_\a^\b $.
\end{proof}

\end{document}
